% BreakpointProfiles.jl: Profile-likelihood uncertainty quantification for dynamical models with fixed changepoints
% Preprint manuscript, PLOS Computational Biology / Bioinformatics style

\documentclass[10pt,letterpaper]{article}
\usepackage[top=0.85in,left=2.75in,footskip=0.75in]{geometry}
\usepackage{amsmath,amssymb}
\usepackage{changepage}
\usepackage{textcomp,marvosym}
\usepackage{cite}
\usepackage{nameref,hyperref}
\usepackage[right]{lineno}
\usepackage[nopatch=eqnum]{microtype}
\DisableLigatures[f]{encoding = *, family = * }
\usepackage[table]{xcolor}
\usepackage{array}
\newcolumntype{+}{!{\vrule width 2pt}}
\newlength\savedwidth
\newcommand\thickcline[1]{%
  \noalign{\global\savedwidth\arrayrulewidth\global\arrayrulewidth 2pt}%
  \cline{#1}%
  \noalign{\vskip\arrayrulewidth}%
  \noalign{\global\arrayrulewidth\savedwidth}%
}
\newcommand\thickhline{\noalign{\global\savedwidth\arrayrulewidth\global\arrayrulewidth 2pt}%
\hline
\noalign{\global\arrayrulewidth\savedwidth}}
\raggedright
\setlength{\parindent}{0.5cm}
\textwidth 5.25in
\textheight 8.75in
\usepackage[aboveskip=1pt,labelfont=bf,labelsep=period,justification=raggedright,singlelinecheck=off]{caption}
\renewcommand{\figurename}{Fig}
\bibliographystyle{plos2025}
\makeatletter
\renewcommand{\@biblabel}[1]{\quad#1.}
\makeatother
\usepackage{lastpage,fancyhdr,graphicx}
\usepackage{epstopdf}
\pagestyle{fancy}
\fancyhf{}
\rfoot{\thepage/\pageref{LastPage}}
\renewcommand{\headrulewidth}{0pt}
\renewcommand{\footrule}{\hrule height 2pt \vspace{2mm}}
\fancyheadoffset[L]{2.25in}
\fancyfootoffset[L]{2.25in}
\lfoot{\today}

\begin{document}
\vspace*{0.2in}

\begin{flushleft}
{\Large
\textbf{BreakpointProfiles.jl: profile-likelihood uncertainty quantification for dynamical models with fixed changepoints}
}
\newline
\\
Mehdi Lotfi\textsuperscript{1*},
\\
\bigskip
\textbf{1} Institute for Medical Microbiology, University Medical Center Greifswald, Greifswald, Germany
\\
\bigskip
\end{flushleft}

\section*{Abstract}
Profile-likelihood estimation (PLE) is a widely used method for assessing the practical identifiability of parameters and for constructing confidence intervals after a model has been fitted to data. For dynamical models whose structure is allowed to change at known or estimated changepoints, uncertainty quantification is complicated by two additional questions: are the estimated segment-specific parameters practically identifiable, and how uncertain are the estimated changepoint locations? We present \texttt{BreakpointProfiles.jl}, a Julia package that provides profile-likelihood-based uncertainty quantification for ordinary differential equation (ODE) models and generic objective functions with fixed changepoints. The package supports conditional parameter profiling (parameters vary, changepoints fixed), joint adaptive profiling (parameters and changepoint locations treated as nuisance variables), and changepoint-location profiling. It implements adaptive D2D-style stepping, multiple optimizer backends, Gaussian and Laplace likelihoods, bootstrap thresholds for non-Gaussian losses, and utilities for summarising, exporting, and plotting results. We demonstrate the package on a COVID-19 SEIRD model with two changepoints identified by the MICA segmentation framework, and we show how joint profiling changes the interpretation of parameter identifiability when changepoint locations are not treated as fixed constants. \texttt{BreakpointProfiles.jl} is designed to be model-agnostic and can be applied to any dynamical model or objective function that exposes a parameter vector and a changepoint vector.

\section*{Author summary}
Many dynamical models are fitted under the assumption that one or more parameters change at specific time points called changepoints. Once the changepoints and the segment-specific parameters have been estimated, it is important to report their uncertainty. We developed a Julia software package, \texttt{BreakpointProfiles.jl}, that automates the computation of profile-likelihood confidence intervals for both the parameters and the changepoint locations. The package can be used with any model that takes a parameter vector and a changepoint vector, and it includes examples for ODE-based epidemiological models. The software fills a gap in the Julia ecosystem for uncertainty quantification in changepoint models and is intended to support reproducible, computationally efficient profiling of complex dynamical systems.

\clearpage
\newgeometry{top=0.85in,left=1in,right=1in,footskip=0.75in}
\linenumbers

\section*{Introduction}
Dynamical systems in biology, epidemiology, and engineering are often described by ordinary differential equations (ODEs) whose parameters may change at discrete time points, called changepoints or breakpoints. Estimating both the changepoint locations and the segment-specific parameters is a well-studied problem, and many algorithms exist for detecting changepoints in time-series data or for fitting piecewise-stationary models \cite{killick2012,truong2020,van_den_burg2020}. However, once a segmented model has been fitted, reporting uncertainty is equally important and is far less standardized.

Profile-likelihood estimation (PLE) is one of the most reliable approaches for practical identifiability analysis and confidence-interval construction in nonlinear models \cite{raue2009,d2d2015}. The idea is simple: for each parameter of interest, one fixes the parameter at a grid of values and re-optimises all remaining parameters. The resulting profile of the likelihood (or loss) function is then compared with a threshold to obtain a confidence interval. Under standard regularity conditions and Gaussian errors, the threshold is given by Wilks' theorem; for other loss functions, bootstrap or other calibration methods can be used.

In changepoint models, three distinct profiling questions arise. First, conditional on the detected changepoints, are the segment-specific parameters practically identifiable? Second, how uncertain are the changepoint locations themselves when all other parameters are re-optimised? Third, if the changepoints are themselves estimated quantities, should they be treated as nuisance variables during parameter profiling, thereby giving joint confidence statements that account for both parameter and changepoint uncertainty? These questions are particularly relevant when the detection algorithm returns a point estimate of the changepoints but does not report their uncertainty.

Several existing tools implement PLE for ODE models. Data2Dynamics (D2D) provides a mature MATLAB-based implementation of profile-likelihood analysis for systems-biology models \cite{d2d2015}, but it does not natively support changepoints. The R package \texttt{PDE} and related frameworks focus on parameter estimation rather than profiling. In the Julia ecosystem, packages such as \texttt{DiffEqParamEstim.jl} and \texttt{StructuralIdentifiability.jl} address parameter estimation and structural identifiability, respectively, but not profile-likelihood confidence intervals for changepoint models. To our knowledge, there is no reusable, open-source package that combines (i) conditional parameter PLE, (ii) changepoint-location PLE, and (iii) joint adaptive profiling with changepoints as nuisance variables, for generic ODE models or objective functions.

We therefore developed \texttt{BreakpointProfiles.jl}, a Julia package that implements these three profiling modes in a model-agnostic way. The package was motivated by the uncertainty-analysis needs of the MICA (Model-Informed Changepoint Analysis) framework \cite{mica}, which detects changepoints in ODE models by comparing the improvement in a loss function against a penalty, but does not automatically report uncertainty. \texttt{BreakpointProfiles.jl} can be used as a standalone tool or as a post-processing step after any changepoint-detection algorithm. It is implemented in pure Julia, supports parallel computation, and can be applied to any model that exposes either an objective function, a simulator, or a plain ODE right-hand side.

\section*{Design and implementation}

\subsection*{Problem statement}
Let $y(t)$ be a multivariate time series observed at integer times $t = 1, \dots, T$. We consider a model that depends on a parameter vector $\theta \in \mathbb{R}^p$ and a vector of $K$ changepoints $\tau = (\tau_1, \dots, \tau_K)$ with $1 < \tau_1 < \dots < \tau_K < T$. Between consecutive changepoints, the model parameters are constant; at a changepoint, some or all parameters may switch to new values. The parameters are partitioned into $n_\text{global}$ global parameters (shared across all segments) and $n_\text{segment}$ segment-specific parameters (one value per segment), so that the total dimension is $p = n_\text{global} + (K+1) \, n_\text{segment}$.

Given a loss function $L(\theta, \tau)$, such as a Gaussian or Laplace negative log-likelihood, the fitted model is obtained by minimising $L$ over $(\theta, \tau)$. In the MICA framework, this minimisation is performed by a genetic algorithm combined with greedy forward-backward segmentation, and the resulting $(\hat{\theta}, \hat{\tau})$ is the point estimate \cite{mica}. \texttt{BreakpointProfiles.jl} takes this point estimate as input and performs profile-likelihood analysis around it.

\subsection*{Three profiling modes}

\paragraph*{Conditional parameter profiling.}
The changepoints are fixed at $\hat{\tau}$ and each scalar parameter $\theta_i$ is profiled by fixing it to a value $\theta_i^*$ and re-optimising all other parameters:
\begin{equation}
\label{eq:conditional}
PL_i(\theta_i^*) = \min_{\theta_{-i}} L(\theta_i^*, \theta_{-i}; \hat{\tau}).
\end{equation}
This is the standard PLE problem and gives the uncertainty of each parameter conditional on the detected changepoints.

\paragraph*{Changepoint-location profiling.}
All parameters are fixed at $\hat{\theta}$ and one changepoint $\tau_k$ is moved over a grid of candidate locations, while the other changepoints remain fixed:
\begin{equation}
\label{eq:cp}
PL_k^{\text{CP}}(\tau_k^*) = L(\hat{\theta}; \tau_1, \dots, \tau_k^*, \dots, \tau_K).
\end{equation}
This profile shows how the loss changes when a single changepoint is perturbed, and it can be used to construct a confidence interval for the changepoint location.

\paragraph*{Joint adaptive profiling.}
The most general case treats the changepoints as nuisance variables during parameter profiling. For each fixed value $\theta_i^*$, both the other parameters and the changepoint locations are jointly optimised:
\begin{equation}
\label{eq:joint}
PL_i^{\text{joint}}(\theta_i^*) = \min_{\theta_{-i}, \tau} L(\theta_i^*, \theta_{-i}; \tau).
\end{equation}
This profile accounts for the fact that the changepoints are themselves estimated, and it can lead to wider confidence intervals or different identifiability conclusions than the conditional profile.

Because joint optimisation over both continuous parameters and discrete changepoint locations is computationally expensive, \texttt{BreakpointProfiles.jl} uses a discrete grid of changepoint candidates with a user-specified window and step size, combined with a fast local optimiser for the continuous parameters at each candidate set. The stepping along the profiled parameter is adaptive, following the D2D-style strategy of expanding the grid until the profile crosses a user-defined threshold \cite{d2d2015}.

\subsection*{Thresholds and likelihoods}
For Gaussian (L2) negative log-likelihoods, Wilks' theorem applies and the 95\% confidence interval is obtained by comparing the profile with the threshold $L(\hat{\theta}, \hat{\tau}) + \chi^2_{1, 0.95}$, where $\chi^2_{1, 0.95} \approx 3.8415$. For Laplace (L1) or other non-Gaussian losses, Wilks' theorem does not hold in general. The package therefore provides a moving-block bootstrap procedure that resamples the residuals around the fitted trajectory to obtain an empirical threshold for the profile-likelihood ratio. The threshold is passed as a user argument, so the package is not tied to a particular error model or confidence level.

\subsection*{Software architecture}
\texttt{BreakpointProfiles.jl} is organised around a single problem type, \texttt{ODEChangepointPLEProblem}, which stores the model (objective, simulator, or ODE function), the data, the loss function, the fitted parameters and changepoints, parameter bounds, and metadata such as the number of global and segment-specific parameters. The package then exposes three high-level functions: \texttt{profile\_parameter}, \texttt{profile\_changepoint}, and \texttt{profile\_parameter\_joint}. Optimiser backends include Evolutionary.jl (genetic algorithm), Optim.jl (LBFGS, Nelder-Mead), NLopt.jl (BOBYQA), and a hybrid global-then-local optimiser. Results are returned as \texttt{ProfileResult} or \texttt{CPProfileResult} objects, which can be summarised, exported to CSV, or plotted.

\section*{Results}

\subsection*{COVID-19 SEIRD example}
We demonstrate the package on the COVID-19 SEIRD model used in the MICA revision. The model contains 8 global parameters (seasonality, transmission, rates of progression through exposed and infectious compartments, and waning immunity) and 8 segment-specific parameters (transition probabilities, reporting probabilities, and vaccination rates). MICA detected two changepoints at days 60 and 150 in a 400-day daily time series from Germany \cite{mica}.

Conditional parameter profiling with the changepoints fixed at [60, 150] identified 17 of 32 parameters as practically identifiable under an L2 Gaussian loss. Several parameters, including probabilities and reporting rates, hit the upper or lower bounds and were classified as non-identifiable. This is consistent with the high dimensionality of the model and the strong correlations among parameters in the later segments.

Changepoint-location profiling with a window of $\pm 30$ days showed that the first changepoint (day 60) is well-localised, with the loss rising sharply when the changepoint is moved away from the estimated location. The second changepoint (day 150) is more uncertain, with a flatter profile, suggesting that the data do not strongly constrain its exact position.

Joint adaptive profiling, in which the changepoints were treated as nuisance variables with a window of $\pm 20$ days and a step of 5 days, produced wider confidence intervals for several parameters and changed the identifiability status of a subset of them. This illustrates that treating changepoints as fixed constants can give an over-optimistic picture of parameter uncertainty.

\subsection*{Illustrative quadratic example}
To validate the implementation against an analytical benchmark, we profiled a simple quadratic loss $L(\theta_1, \theta_2) = (\theta_1 - 3)^2 + (\theta_2 - 5)^2$. The conditional profile for $\theta_1$ is $PL_1(\theta_1^*) = (\theta_1^* - 3)^2$, and the exact 95\% confidence interval is $[3 - \sqrt{3.84}, 3 + \sqrt{3.84}] \approx [1.04, 4.96]$. The package recovered this interval to within the requested numerical tolerance, confirming the correctness of the adaptive stepping and bound handling.

\section*{Discussion}
\texttt{BreakpointProfiles.jl} provides a missing piece of infrastructure for uncertainty quantification in changepoint models of dynamical systems. By separating the tasks of changepoint detection and profile-likelihood analysis, it allows researchers to use their preferred detection algorithm while still obtaining rigorous uncertainty estimates. The joint profiling mode is particularly important because it avoids the common but questionable practice of treating estimated changepoints as known constants.

The package has several limitations. First, it assumes that the number of changepoints $K$ is fixed; it does not perform model selection over $K$. Second, the joint profiling grid is exhaustive over changepoint candidates and therefore scales combinatorially with $K$. For models with more than a few changepoints, the user should restrict the window and step size or use a more efficient sampling strategy. Third, the bootstrap threshold for L1 losses is computationally expensive and should be used with care; we recommend it primarily as a sensitivity check rather than as the default.

Future work includes extending the package to support stochastic differential equations, implicit likelihoods (e.g. via simulation-based inference), and more efficient joint optimisation over changepoint configurations. We also plan to register the package in the Julia General registry and to provide additional application examples beyond epidemiology.

\section*{Availability and implementation}
\texttt{BreakpointProfiles.jl} is written in Julia and is available at the project repository. The package is released under the MIT license. Documentation and examples, including the COVID-19 SEIRD workflow and the quadratic sanity test, are included in the repository. The software is currently under development; a registered version will be released after the peer review of this manuscript.

\section*{Acknowledgments}
The package was developed as part of the MICA revision. The author thanks the reviewers for constructive comments that motivated the implementation of joint profiling and changepoint-location confidence intervals.

\nolinenumbers

\bibliography{breakpoint_profiles}

\end{document}
